% 原书物理页=印刷页49--52；第 2 章章首（导语）与 2.1 全节.
\chapter{Lie 代数}\label{chap:lie-2}

本章致力于研究 Lie 代数的结构及其对表示理论的诸般影响. 这里引入的绝大多数概念在任意基域 $\mathbb K$ 上都成立.

\section{Lie 代数的泛包络代数}\label{sec:lie-2-1}

回顾一下，任何 结合代数都可以视为一个 $\mathbb K$-Lie 代数，其 Lie 括号为 $[a,b]=ab-ba$.

\begin{theorem}\label{thm:lie-2-1-1}
设 $\mathfrak g$ 是一个 $\mathbb K$-Lie 代数. 则存在一个偶对 $(\mathcal U\mathfrak g,\iota)$，由
\begin{itemize}[label=--]
\item 一个 $\mathbb K$-结合代数 $\mathcal U\mathfrak g$ 与
\item 一个 Lie 代数态射 $\mathfrak g\longrightarrow\mathcal U\mathfrak g$
\end{itemize}
组成，且满足下述\newterm{lie-universal-property}{泛性质}{universal property}：对任何 $\mathbb K$-结合代数 $A$，映射
\[
\operatorname{Hom}_{\mathbb K-\mathrm{alg}}(\mathcal U\mathfrak g,A)\longrightarrow\operatorname{Hom}_{\mathbb K-\mathrm{ALie}}(\mathfrak g,A),\qquad f\mapsto f\circ\iota
\]
是双射.
\end{theorem}

泛性质保证了偶对 $(\mathcal U\mathfrak g,\iota)$ 在相差唯一同构意义下的唯一性. 对于喜欢范畴论的读者，可以把该定理改述成下面两种说法之一：
\begin{itemize}[label=--]
\item 对任何 $\mathfrak g$，函子 $A\mapsto\operatorname{Hom}_{\mathbb K-\mathrm{ALie}}(\mathfrak g,A)$ 是\newterm{lie-representable-functor}{可表的}{representable}.
\item 把一个结合代数映到其底 Lie 代数的函子具有\newterm{lie-left-adjoint}{左伴随}{left adjoint}. 这种表述看似更强，因为它蕴含了 $\mathfrak g\mapsto\mathcal U\mathfrak g$ 的函子性，但实际上二者等价.
\end{itemize}

\begin{proof}
考虑 $\mathsf T\mathfrak g=\bigoplus_{n\in\mathbb N}\mathsf T^n\mathfrak g$，其中 $\mathsf T^n\mathfrak g=\mathfrak g^{\otimes n}=\mathfrak g\otimes_{\mathbb K}\mathfrak g\otimes_{\mathbb K}\cdots\otimes_{\mathbb K}\mathfrak g$. 在这个 $\mathbb K$-向量空间上赋予拼接乘积，它把 $\mathsf T^n\mathfrak g\times\mathsf T^m\mathfrak g$ 映入 $\mathsf T^{n+m}\mathfrak g$，并且在基本张量上由
\[
(x_1\otimes\cdots\otimes x_n)(y_1\otimes\cdots\otimes y_m)=x_1\otimes\cdots\otimes x_n\otimes y_1\otimes\cdots\otimes y_m
\]
给出. 于是我们得到一个\newterm{lie-graded}{分次}{graded}的 $\mathbb K$-结合代数 $\mathsf T\mathfrak g$，称为 $\mathfrak g$ 的\newterm{lie-tensor-algebra}{张量代数}{tensor algebra}.
注意这一构造实际上只依赖于 $\mathfrak g$ 的底 $\mathbb K$-向量空间，因而适用于任意 $\mathbb K$-向量空间. 代数 $\mathsf T\mathfrak g$ 配备映射 $\mathfrak g=\mathsf T^1\mathfrak g\hookrightarrow\mathsf T\mathfrak g$ 后具有如下泛性质：对任何 $\mathbb K$-结合代数 $A$，映射 $\operatorname{Hom}_{\mathbb K-\mathrm{alg}}(\mathsf T\mathfrak g,A)\longrightarrow\operatorname{Hom}_{\mathbb K}(\mathfrak g,A)$ 是双射. 换言之，函子 $V\mapsto\mathsf TV$ 是把代数映到其底向量空间的函子的左伴随. 这一验证很容易，但我们不给出细节，可参看任何代数教材（例如 Lang 的书）.
现在考虑由形如
\[
x\otimes y-y\otimes x-[x,y]\qquad\in\mathsf T^{\le 2}\mathfrak g
\]
的元素生成的\newterm{lie-two-sided-ideal}{双边理想}{two-sided ideal} $I$. 注意它不是一个齐次理想. 最后令 $\mathcal U\mathfrak g:=\mathsf T\mathfrak g/I$. 这是一个结合代数，不分次，但仍带有\newterm{lie-filtration}{滤过}{filtration}：若 $\mathcal U\mathfrak g^{\le n}$ 表示 $\mathsf T^{\le n}\mathfrak g$ 的像，则 $\mathcal U\mathfrak g^{\le n}\,\mathcal U\mathfrak g^{\le m}\subset\mathcal U\mathfrak g^{\le n+m}$. 由构造，有
\[
\operatorname{Hom}_{\mathbb K-\mathrm{alg}}(\mathcal U\mathfrak g,A)=\{f\in\operatorname{Hom}_{\mathbb K-\mathrm{alg}}(\mathsf T\mathfrak g,A),\ \forall x,y\in\mathfrak g,\ f(x)f(y)-f(y)f(x)=f([x,y])\}.
\]
于是针对配备 $\iota:\mathfrak g\hookrightarrow\mathsf T\mathfrak g\twoheadrightarrow\mathcal U\mathfrak g$ 的 $\mathcal U\mathfrak g$ 所宣布的泛性质，便由 $\mathsf T\mathfrak g$ 的泛性质推出.
\end{proof}

\noindent\textbf{应用.}\quad 设 $(V,\rho)$ 是 $\mathfrak g$ 的一个 $\mathbb K$-表示（也称 $V$ 是一个左 \newterm{lie-g-module}{$\mathfrak g$-模}{$\mathfrak g$-module}），则 Lie 代数态射 $\rho:\mathfrak g\longrightarrow\mathfrak{gl}(V)$ 由一个（唯一的）代数态射 $\mathcal U\mathfrak g\longrightarrow\operatorname{End}_{\mathbb K}(V)$ 导出，后者使 $V$ 成为左 $\mathcal U\mathfrak g$-模. 于是我们得到一个范畴等价
\[
\operatorname{Rep}_{\mathbb K}(\mathfrak g)=\mathfrak g\text{--Mod}\ \xrightarrow{\ \sim\ }\ \mathcal U\mathfrak g\text{--Mod}
\]
连接左 $\mathfrak g$-模范畴与左 $\mathcal U\mathfrak g$-模范畴. 特别地，$\mathfrak g$ 的不可约表示对应于 $\mathcal U\mathfrak g$ 的单模.

\noindent\textbf{例.}\quad 首先注意 $\mathcal U\mathfrak{gl}_n(\mathbb K)\neq\mathrm M_n(\mathbb K)$！！！例如取 $n=2$：元素 $E=E_{12}$ 在 $\mathrm M_2(\mathbb K)$ 中是 $2$ 阶幂零的. 但在前面构造的表示 $V_m$ 中，它的像是 $m+1$ 阶幂零的；由于 $m$ 任意，这说明在 $\mathcal U\mathfrak{gl}_2$ 中它的像根本不是幂零的！事实上，$\mathcal U\mathfrak g$ 一般是无限维的.
\begin{enumerate}[label=\roman*)]
\item $\mathfrak g=\mathbb K$. 可验证 $\mathcal U\mathfrak g=\mathbb K[T]$，其中 $\mathfrak g=\mathbb K\xrightarrow{\ \sim\ }\mathbb K.T$ 即符合要求.
\item $\mathfrak g$ 交换. 此时对 $\mathsf T\mathfrak g$ 作商所用的理想由 $x\otimes y-y\otimes x$ 生成. 所得代数是\newterm{lie-symmetric-algebra}{对称代数}{symmetric algebra} $S\mathfrak g=\bigoplus_{n\in\mathbb N}S^n\mathfrak g$，它交换且分次，其构造适用于任意 $\mathbb K$-向量空间. 函子 $V\mapsto SV$ 是把交换代数映到其底向量空间的函子的左伴随. 我们有 $S^n\mathfrak g=(\mathsf T^n\mathfrak g)^{\mathfrak S_n}$（在对称群 $\mathfrak S_n$ 作用下的\newterm{lie-coinvariant}{余不变量}{coinvariants}）. $\mathfrak g$ 的任一组基的选取诱导同构 $S\mathfrak g\xrightarrow{\ \sim\ }\mathbb K[T_1,\cdots,T_{\dim\mathfrak g}]$.
\item $\mathfrak g=\mathfrak{sl}_2(\mathbb K)$. 可以给出如下\newterm{lie-presentation}{表现}{presentation}：$\mathcal U\mathfrak{sl}_2$ 是由生成元 $e,f,h$ 按关系 $he-eh=2e$、$hf-fh=-2f$ 与 $ef-fe=h$ 给出的结合代数. 它含有一个引人注目的元素，即由
\[
C:=ef+fe+\tfrac12 h^2.
\]
给出的\newterm{lie-casimir-operator}{Casimir 算子}{Casimir operator}.
\begin{exercise*}
\begin{enumerate}[label=\alph*)]
\item 证明 $C$ 与 $e,f,h$ 交换，从而属于 $\mathcal U\mathfrak g$ 的中心.
\item 若 $\mathbb K=\mathbb C$（更一般地，若 $\mathbb K$ 代数闭），由此推出 $C$ 在每个不可约表示上以一个纯量作用.
\item 证明在表示 $V_m$ 上，这个纯量为 $\lambda_m=\tfrac12 m(m+2)$.
\end{enumerate}
\end{exercise*}
\item $\mathfrak g=\mathfrak{so}_3(\mathbb R)$. 可以给出如下表现：$\mathcal U\mathfrak{so}_3(\mathbb R)$ 是由生成元 $J_x,J_y,J_z$ 按关系 $J_xJ_y-J_yJ_x=J_z$、$J_yJ_z-J_zJ_y=J_x$ 与 $J_zJ_x-J_xJ_z=J_y$ 给出的结合代数. 考虑算子
\[
C=J_x^2+J_y^2+J_z^2.
\]
\begin{exercise*}
\begin{enumerate}[label=\alph*)]
\item 证明 $C$ 属于 $\mathcal U\mathfrak{so}_3(\mathbb R)$ 的中心.
\item 验证 $C$ 在 $V_m$ 上以 $-\tfrac14 m(m+2)$ 相乘.
\end{enumerate}
\end{exercise*}
\end{enumerate}

\begin{numpara}[泛包络代数与微分算子]\label{np:lie-2-1-2}
设一个 Lie 群光滑地作用在一个流形 $M$ 上. 前面已经看到如何把这个作用求导，得到 $\operatorname{Lie}(G)$ 在 $\mathcal C^\infty(M)$ 上由导子给出的作用，它由一个 Lie 代数（反）态射
\[
d\alpha:\operatorname{Lie}(G)\longrightarrow\mathcal V(M)=\operatorname{Der}(\mathcal C^\infty(M)).
\]
给出. 设 $\operatorname{Diff}(\mathcal C^\infty(M))$ 是 $M$ 的\newterm{lie-differential-operator}{微分算子}{differential operator}代数，即 $\operatorname{End}_{\mathbb R}(\mathcal C^\infty(G))$ 中由导子生成的子 $\mathbb R$-代数.\footnote{译注：原文如此，其中 $\mathcal C^\infty(G)$ 疑为 $\mathcal C^\infty(M)$ 之误.} $\mathcal U\operatorname{Lie}(G)$ 的泛性质使我们能把 $d\alpha$ 扩张成一个（反）代数态射
\[
d\alpha:\mathcal U\operatorname{Lie}(G)\longrightarrow\operatorname{Diff}(\mathcal C^\infty(M))
\]
从\newterm{lie-universal-enveloping-algebra}{泛包络代数}{universal enveloping algebra}到 $M$ 的微分算子代数.
\end{numpara}

\begin{example*}
（$\mathrm{SL}_3(\mathbb R)$ 在球面 $S^2$ 上的作用与球面 Laplace 算子的谱）此情形下我们已经算得
\[
d\alpha(J_z)=x\frac{\partial}{\partial y}-y\frac{\partial}{\partial x},\qquad
d\alpha(J_y)=z\frac{\partial}{\partial x}-x\frac{\partial}{\partial z},\qquad
d\alpha(J_x)=y\frac{\partial}{\partial z}-z\frac{\partial}{\partial y}.
\]
利用 $x^2+y^2+z^2=1$，可得
\[
d\alpha(C)=d\alpha\bigl(J_x^2+J_y^2+J_z^2\bigr)
=\Delta-\Bigl(x\frac{\partial}{\partial x}+y\frac{\partial}{\partial y}+z\frac{\partial}{\partial z}\Bigr)^{\!2}
-\Bigl(x\frac{\partial}{\partial x}+y\frac{\partial}{\partial y}+z\frac{\partial}{\partial z}\Bigr),
\]
其中 $\Delta=\frac{\partial^2}{\partial x^2}+\frac{\partial^2}{\partial y^2}+\frac{\partial^2}{\partial z^2}$ 是 $\mathbb R^3$ 的 Laplace 算子. 引入球面 Laplace 算子
\[
\Delta^{sph}:\ \begin{array}[t]{rcl}
\mathcal C^\infty(S^2)&\to&\mathcal C^\infty(S^2)\\
f&\mapsto&\Delta(\tilde f)|_{S^2}
\end{array}
\]
其中 $\tilde f$ 表示 $f$ 到 $\mathbb R^3\setminus\{0\}$ 的延拓，由 $\tilde f(v):=f(v/\|v\|)$ 给出. 对 $\lambda\mapsto\tilde f(\lambda v)$ 求导，可知 $\bigl(x\frac{\partial}{\partial x}+y\frac{\partial}{\partial y}+z\frac{\partial}{\partial z}\bigr)(\tilde f)=0$，最终得到
\[
d\alpha(C)=\Delta^{sph}.
\]
由此推出 $\Delta^{sph}$ 在 $\mathcal C^\infty(S^2)$ 上的谱：
\[
\operatorname{Sp}(\Delta^{sph})=\{-k(k+1),\ k\in\mathbb N\}.
\]
事实上，我们已经看到，由 $x,y,z$ 的多项式函数组成的子空间 $\mathcal P(S^2)$ 分解为直和 $\bigoplus_{k\in\mathbb N}V_{2k}$，其中 $V_{2k}$ 表示 $\mathrm{SO}_3(\mathbb R)$ 的维数为 $2k+1$ 的不可约表示，且 $C$ 在 $V_{2k}$ 上以纯量 $-k(k+1)$ 作用. 这给出了包含关系 $\operatorname{Sp}(\Delta^{sph})\supset\{-k(k+1),\ k\in\mathbb N\}$. 为证另一侧的包含，把 $\mathcal C^\infty(S^2)$ 嵌入 $L^2(S^2)$. $L^2(S^2)$ 的 Hermite 内积在 $\mathrm{SO}_3(\mathbb R)$ 作用下不变，因此算子 $d\alpha(C)$ 是自伴的. 此外，和 $\bigoplus_{k\in\mathbb N}V_{2k}$ 是正交的，并且由多项式函数的稠密性，我们有 $L^2(S^2)=\widehat{\bigoplus}_{k\in\mathbb N}V_{2k}$（Hilbert 和）. 于是 $\operatorname{Sp}(\Delta^{sph})\subset\{-k(k+1),\ k\in\mathbb N\}$ 得证.
\end{example*}

\begin{remark*}
在 $G$ 通过\emph{右}平移作用于 $M=G$ 的特殊情形，代数态射 $\mathcal U\operatorname{Lie}(G)\longrightarrow\operatorname{Diff}(\mathcal C^\infty(G))$ 的像包含在\emph{左}不变微分算子子代数 $\operatorname{Diff}(\mathcal C^\infty(G))^L$ 中. 可以证明这个态射是同构
\[
\mathcal U\operatorname{Lie}(G)\xrightarrow{\ \sim\ }\operatorname{Diff}(\mathcal C^\infty(G))^L.
\]
\end{remark*}

\begin{numpara}[Poincaré--Birkhoff--Witt 定理]\label{np:lie-2-1-3}
这里设 $\mathfrak g$ 在 $\mathbb K$ 上维数有限.
\begin{theorem*}[第一种表述]
设 $x_1,\cdots,x_n$ 是 $\mathfrak g$ 在 $\mathbb K$ 上的一组基，则形如 $x_{i_1}x_{i_2}\cdots x_{i_k}$（其中 $i_1\le i_2\le\cdots\le i_k$）的单子构成 $\mathcal U\mathfrak g$ 在 $\mathbb K$ 上的一组基.
\end{theorem*}
我们在此不证明该定理（例如可参看 Carter 的书）. 不过要指出，容易看出这族元素是生成元：只需迭代形如 $x_ix_j=x_jx_i+[x_i,x_j]$（当 $i>j$ 时）的关系即可（做练习）. 而证明这族元素线性无关则要繁琐得多.
现在给出同一结果的一个更内蕴的表述. 为此，回顾 $\mathcal U\mathfrak g$ 从张量代数 $\mathsf T\mathfrak g$ 的同一滤过继承来的递增滤过 $(\mathcal U\mathfrak g^{\le k})_{k\in\mathbb N}$（见上文存在性定理的证明）. 记 $\operatorname{gr}(\mathcal U\mathfrak g):=\bigoplus_{k\in\mathbb N}\mathcal U\mathfrak g^{\le k}/\mathcal U\mathfrak g^{\le k-1}$ 为相配的分次代数. 容易验证它是交换的，因为两个长度为 $k$ 的单子的换位子是长度 $\le k-1$ 的单子的线性组合. 于是，典范态射 $\mathsf T\mathfrak g=\operatorname{gr}(\mathsf T\mathfrak g)\longrightarrow\operatorname{gr}(\mathcal U\mathfrak g)$ 可经由对称代数（见上文例 ii)）分解，从而给出一个分次代数态射 $S\mathfrak g\longrightarrow\operatorname{gr}(\mathcal U\mathfrak g)$.
\begin{theorem*}[第二种表述]
上述态射 $S\mathfrak g\longrightarrow\operatorname{gr}(\mathcal U\mathfrak g)$ 是同构.
\end{theorem*}
无论采用哪种表述，都得到以下结论：
\begin{corollary*}
\begin{enumerate}[label=\roman*)]
\item 映射 $\iota:\mathfrak g\longrightarrow\mathcal U\mathfrak g$ 是单射.
\item 对任何 Lie 子代数 $\mathfrak h\subset\mathfrak g$，$\mathcal U\mathfrak h$ 可等同于 $\mathcal U\mathfrak g$ 的一个子代数.
\end{enumerate}
\end{corollary*}
注意 i) 也可由 Ado 定理推出，该定理说 $\mathfrak g$ 可嵌入某个 $\mathfrak{gl}(V)$，从而嵌入 $\operatorname{End}_{\mathbb K}(V)$.
\end{numpara}

\begin{numpara}[双代数结构]\label{np:lie-2-1-4}
以下面的注记作结，它表明泛包络代数拥有比结合代数更多的结构. 考虑映射 $\mathfrak g\longrightarrow\mathcal U\mathfrak g\otimes_{\mathbb K}\mathcal U\mathfrak g$，它把 $x$ 映到 $\iota(x)\otimes 1+1\otimes\iota(x)$. 可以证明这是一个 Lie 代数态射（习题）. 因此它诱导一个代数态射 $\Delta:\mathcal U\mathfrak g\longrightarrow\mathcal U\mathfrak g\otimes_{\mathbb K}\mathcal U\mathfrak g$. 进而可以证明等式 $(\Delta\otimes\mathrm{id})\circ\Delta=(\mathrm{id}\otimes\Delta)\circ\Delta$（\newterm{lie-coassociativity}{余结合性}{coassociativity}公理）. 我们说 $\Delta$ 是一个\newterm{lie-coproduct}{余积}{coproduct}，并说 $\mathcal U\mathfrak g$ 是一个\newterm{lie-bialgebra}{双代数}{bialgebra}（甚至是 \newterm{lie-hopf-algebra}{Hopf 代数}{Hopf algebra}）. 我们在此不再深入，但这一观点在\newterm{lie-quantum-group}{量子群}{quantum groups}理论中是根本性的. 量子群并不是真正的群，而更确切地说是经典泛包络代数的``形变''.
\end{numpara}
